\documentclass[11pt,letterpaper]{article}

% Essential Packages
\usepackage[margin=1in]{geometry}
\usepackage{amsmath,amssymb}
\usepackage{graphicx}
\usepackage{circuitikz}
\usepackage{tikz}
\usetikzlibrary{shapes,arrows,positioning,calc}
\usepackage{booktabs}
\usepackage{multirow}
\usepackage{xcolor}
\usepackage{fancyhdr}
\usepackage{tcolorbox}
\usepackage{enumitem}
\usepackage{float}
\usepackage{listings}
\usepackage{hyperref}

% Colors
\definecolor{nlcblue}{RGB}{0,51,102}
\definecolor{alertred}{RGB}{204,0,0}
\definecolor{safgreen}{RGB}{0,128,0}

% Header/Footer
\pagestyle{fancy}
\fancyhead[L]{\textbf{NLC STEM Club UAV Project}}
\fancyhead[R]{\textbf{Subsystem Assignments}}
\fancyfoot[C]{\thepage}

% Title
\title{\Huge\textbf{Subsystem Assignments}\\
\Large Team Organization and Responsibilities\\
\large NLC STEM Club - Electrical Systems Team}
\author{February 2026}
\date{}

\begin{document}

\maketitle
\thispagestyle{empty}
\newpage
\section{SUBSYSTEM ASSIGNMENT \& TRACKING DOCUMENT}

\begin{itemize}
  \item Fixed-Wing Environmental Monitoring UAV
  \item NLC STEM Club - Electrical Systems Team
\end{itemize}

\section{SUBSYSTEM 1: POWER TEAM}

\section{TEAM MEMBERS}

Owner:         \underline{\hspace{4.8cm}}  Phone: \underline{\hspace{1.7999999999999998cm}}  Email: \underline{\hspace{1.95cm}}\par

Assistant 1:   \underline{\hspace{4.8cm}}  Phone: \underline{\hspace{1.7999999999999998cm}}  Email: \underline{\hspace{1.95cm}}\par

Assistant 2:   \underline{\hspace{4.8cm}}  Phone: \underline{\hspace{1.7999999999999998cm}}  Email: \underline{\hspace{1.95cm}}\par

\section{SUBSYSTEM SCOPE}

You are responsible for ALL power generation, distribution, and safety for the aircraft.\par

This includes battery management, voltage regulation, and ensuring stable power to all\par

flight-critical and research payload systems.\par

\section{MISSION-CRITICAL IMPORTANCE:}

Without stable, properly distributed power, nothing else works. You are the foundation.\par

\section{COMPONENTS ASSIGNED TO YOU}

\begin{itemize}
  \item OVONIC 4S LiPo Battery (14.8V, 2200mAh)
  \item B6 Battery Charger (80W)
  \item LiPo Safety Bag
  \item 40A ESC (includes 5V 3A BEC)
  \item External 5V UBEC (5-6A rating) [TO BE ORDERED]
  \item XT60 Connectors (Y-harness or splitter)
  \item Wire (12-14 AWG for high current, 18-20 AWG for BEC outputs)
  \item Heat shrink tubing (various sizes)
  \item Multimeter
\end{itemize}

WEEK 1 DELIVERABLES (Days 1-7)\par

\begin{itemize}
  \item Battery Management Setup
  \item Charge battery to full (16.8V or 4.2V per cell) using B6 charger
  \item Verify balance: All 4 cells within 0.05V of each other
  \item Document battery voltage after full charge: \_\_\_\_\_\_V
  \item Label battery with charge date and voltage
  \item Create battery log sheet (date, voltage before use, voltage after use, flight time)
\end{itemize}

\begin{itemize}
  \item ESC BEC Validation
  \item Connect battery to ESC (XT60)
  \item Measure ESC BEC output with multimeter
  \item Expected: 5.0V $\pm$ 0.1V    Actual: \_\_\_\_\_\_V    Status: PASS / FAIL
  \item Photograph ESC BEC output measurement
\end{itemize}

\begin{itemize}
  \item External UBEC Setup (when arrives)
  \item Verify UBEC received and undamaged
  \item Solder input wires (14.8V from battery) - use 16-18 AWG wire
  \item Solder output wires (5V output) - use 18-20 AWG wire
  \item Apply heat shrink to all solder joints
  \item Connect UBEC input to battery (via test leads, not permanent yet)
  \item Measure UBEC output with multimeter
  \item Expected: 5.0V $\pm$ 0.1V    Actual: \_\_\_\_\_\_V    Status: PASS / FAIL
  \item Photograph UBEC and voltage measurement
\end{itemize}

\begin{itemize}
  \item Power Distribution Planning
  \item Design XT60 Y-harness or splitter:
  \item Battery (+) $\rightarrow$ ESC (+) AND External UBEC (+)
  \item Battery (-) $\rightarrow$ ESC (-) AND External UBEC (-)
  \item Calculate wire lengths needed (measure from battery compartment to components)
  \item Sketch wiring layout on paper (including wire routing)
  \item Order/purchase any missing connectors or wire
\end{itemize}

WEEK 2 DELIVERABLES (Days 8-14)\par

\begin{itemize}
  \item Power Distribution Assembly
  \item Build XT60 Y-harness (or install splitter)
  \item Solder all connections, apply heat shrink
  \item Label wires clearly: "Battery Main", "To ESC", "To UBEC"
  \item Test continuity with multimeter (0$\Omega$ from battery + to ESC + and UBEC +)
\end{itemize}

\begin{itemize}
  \item Rail 1 (Flight-Critical) Assembly
  \item Create power distribution harness for Rail 1:
  \item ESC BEC (+5V) $\rightarrow$ FC, GPS, Receiver, Arduino
  \item Use 20-22 AWG wire, red for +5V, black for GND
  \item Create common ground bus (all grounds connected)
  \item Heat shrink all connections
  \item Label: "RAIL 1 - Flight-Critical"
\end{itemize}

\begin{itemize}
  \item Rail 2 (Servo Power) Assembly
  \item Create power distribution harness for Rail 2:
  \item External UBEC (+5V) $\rightarrow$ All 4 servo red wires
  \item Use 18-20 AWG wire (servos draw more current)
  \item Create servo power bus (all 4 servo red wires to one +5V source)
  \item Create servo ground bus (all 4 servo brown wires to one GND)
  \item Heat shrink all connections
  \item Label: "RAIL 2 - Servo Power"
\end{itemize}

\begin{itemize}
  \item Full System Power Test
  \item Connect full power system:
  \item Battery $\rightarrow$ Y-harness $\rightarrow$ ESC + UBEC
  \item ESC BEC $\rightarrow$ Rail 1
  \item External UBEC $\rightarrow$ Rail 2
  \item Power on system (no motor/servos connected yet)
  \item Measure voltages:
  \item Battery voltage: \_\_\_\_\_\_V (expect 14-16.8V)
  \item Rail 1 voltage: \_\_\_\_\_\_V (expect 5.0 $\pm$ 0.1V)
  \item Rail 2 voltage: \_\_\_\_\_\_V (expect 5.0 $\pm$ 0.1V)
  \item Verify common ground: Measure resistance from battery (-) to Rail 1 GND and
  \item Rail 2 GND. Should be 0$\Omega$ (continuity beep).
  \item Photograph entire power distribution setup
\end{itemize}

\begin{itemize}
  \item Current Measurements (coordinate with Integration Team)
  \item Measure Rail 1 idle current (only FC + GPS + RX powered): \_\_\_\_\_\_mA
  \item Expected: \textasciitilde{}320-400mA
  \item Measure Rail 1 with Arduino active: \_\_\_\_\_\_mA
  \item Expected: \textasciitilde{}660mA
  \item Measure Rail 2 idle current (servos idle): \_\_\_\_\_\_mA
  \item Expected: \textasciitilde{}40mA
  \item Measure Rail 2 with 2 servos moving: \_\_\_\_\_\_mA
  \item Expected: \textasciitilde{}400-500mA
  \item Measure Rail 2 with all 4 servos moving: \_\_\_\_\_\_mA
  \item Expected: \textasciitilde{}720-1200mA
\end{itemize}

\begin{itemize}
  \item Voltage Sag Test (CRITICAL)
  \item Connect multimeter probes to Flight Controller 5V input (leave connected)
  \item Power system, note idle voltage: \_\_\_\_\_\_V
  \item Command all 4 servos to move simultaneously (full deflection)
  \item Note voltage during servo movement: \_\_\_\_\_\_V
  \item PASS CRITERIA: Voltage must stay >4.8V, ideally >4.9V
  \item If voltage drops below 4.8V $\rightarrow$ FAILURE $\rightarrow$ Increase BEC capacity or fix wiring
  \item Status: PASS / FAIL
\end{itemize}

\begin{itemize}
  \item Documentation
  \item Create power system diagram with actual wire colors and lengths
  \item Take photos of all major connections
  \item Write 1-page summary: "Power System Integration Report"
  \item What you built
  \item Test results
  \item Issues encountered and how you solved them
  \item Lessons learned
\end{itemize}

\section{SKILLS YOU WILL DEVELOP}

\begin{itemize}
  \item LiPo battery management and charging
  \item Soldering high-current connectors (XT60)
  \item Voltage regulation and power distribution
  \item Multimeter proficiency (voltage, current, continuity)
  \item Wire gauge selection for different current loads
  \item Heat shrink application for professional finish
  \item Systematic testing and validation procedures
\end{itemize}

\section{POTENTIAL ISSUES \& TROUBLESHOOTING}

Problem: Battery won't charge or charger shows error\par

\begin{itemize}
  \item $\rightarrow$ Check balance connector (all 5 pins connected)
  \item $\rightarrow$ Verify cell voltages (all cells >3.0V, <4.3V)
  \item $\rightarrow$ Ensure charger set to LiPo mode, 4S
\end{itemize}

Problem: BEC output voltage is wrong (not 5.0V)\par

\begin{itemize}
  \item $\rightarrow$ Verify battery voltage is adequate (>13.5V)
  \item $\rightarrow$ Check BEC solder connections (cold joint?)
  \item $\rightarrow$ Test BEC with minimal load (just multimeter, no components)
  \item $\rightarrow$ If still wrong, BEC may be defective
\end{itemize}

Problem: Voltage sag test fails (voltage drops below 4.8V)\par

\begin{itemize}
  \item $\rightarrow$ Check wire gauge (may be too thin, causing voltage drop)
  \item $\rightarrow$ Check solder joints (high resistance = voltage drop)
  \item $\rightarrow$ Verify External UBEC is providing power to servos (not ESC BEC)
  \item $\rightarrow$ Measure current draw during test (may exceed BEC capacity)
\end{itemize}

\section{NOTES \& QUESTIONS}

[Space for team to write notes during work]\par

\section{SUBSYSTEM 2: FLIGHT ELECTRONICS TEAM}

\section{TEAM MEMBERS}

Owner:         \underline{\hspace{4.8cm}}  Phone: \underline{\hspace{1.7999999999999998cm}}  Email: \underline{\hspace{1.95cm}}\par

Assistant 1:   \underline{\hspace{4.8cm}}  Phone: \underline{\hspace{1.7999999999999998cm}}  Email: \underline{\hspace{1.95cm}}\par

Assistant 2:   \underline{\hspace{4.8cm}}  Phone: \underline{\hspace{1.7999999999999998cm}}  Email: \underline{\hspace{1.95cm}}\par

\section{SUBSYSTEM SCOPE}

You are responsible for the BRAIN and CONTROL of the aircraft. This includes flight\par

controller configuration, GPS navigation, radio control, and servo actuation.\par

\section{MISSION-CRITICAL IMPORTANCE:}

You control the aircraft's behavior. If you configure incorrectly, we crash.\par

Precision and attention to detail are essential.\par

\section{COMPONENTS ASSIGNED TO YOU}

\begin{itemize}
  \item MATEK F405-WMN Flight Controller
  \item BN-220 GPS Module (with compass)
  \item HAPPYMODEL ELRS EP2 Receiver
  \item Jumper Smart ELRS Transmitter (ground station)
  \item 4$\times$ RGBZONE ES08MA II Metal Gear Servos
  \item 40A ESC (for motor control signal)
  \item FLASH HOBBY D2830EVO Motor (1000KV)
  \item Folding Propeller (1060 or 1160) - DO NOT INSTALL YET
  \item Laptop with USB cable (for FC configuration)
\end{itemize}

WEEK 1 DELIVERABLES (Days 1-7)\par

\begin{itemize}
  \item Flight Controller Setup
  \item Install flight controller configuration software on laptop:
  \item Option A: iNav Configurator (recommended for fixed-wing)
  \item Option B: ArduPilot Mission Planner
  \item Team Decision: We will use: \underline{\hspace{2.25cm}}
  \item Connect FC to laptop via USB
  \item Verify FC powers on (LED should light up)
  \item Connect to FC in configuration software
  \item Check current firmware version: \underline{\hspace{2.25cm}}
  \item Flash latest stable firmware if needed
  \item iNav: Version 7.0 or later
  \item ArduPilot: Plane version 4.5 or later
  \item Screenshot of FC connected in software
\end{itemize}

\begin{itemize}
  \item GPS Module Setup
  \item Inspect GPS module and antenna
  \item Identify GPS connector type on FC (likely JST-SH 4-pin)
  \item Wire GPS to FC UART1:
  \item GPS VCC (red) $\rightarrow$ FC 5V or dedicated GPS 5V pad
  \item GPS TX (yellow) $\rightarrow$ FC RX1
  \item GPS RX (green) $\rightarrow$ FC TX1
  \item GPS GND (black) $\rightarrow$ FC GND
  \item In FC software, enable GPS on UART1
  \item Set GPS protocol: AUTO or UBX
  \item Set baud rate: 9600 or 115200 (try both if one doesn't work)
  \item Power FC via USB, take GPS module OUTDOORS
  \item Wait 5 minutes for GPS lock
  \item Verify GPS working:
  \item Satellite count: \_\_\_\_\_\_ (expect >8 for good lock)
  \item Latitude: \underline{\hspace{2.1cm}}  Longitude: \underline{\hspace{2.1cm}}
  \item GPS coordinates match your location? YES / NO
  \item Screenshot of GPS locked in software
\end{itemize}

\begin{itemize}
  \item Receiver Binding \& Configuration
  \item Read ELRS EP2 receiver binding instructions
  \item Power receiver (5V + GND from bench power supply or FC)
  \item Put receiver in bind mode (press bind button during power-up OR use WiFi method)
  \item Put transmitter in bind mode (follow Jumper TX instructions)
  \item Verify bind successful (receiver LED behavior changes)
  \item Wire receiver to FC:
  \item RX VCC (red) $\rightarrow$ FC 5V
  \item RX GND (black) $\rightarrow$ FC GND
  \item RX TX (white/yellow) $\rightarrow$ FC RX2 (or dedicated SBUS pad)
  \item In FC software, enable CRSF or SBUS protocol on UART2 (check which receiver uses)
  \item Set baud rate: 420000 (CRSF) or 100000 (SBUS)
  \item Verify receiver channels active:
  \item Move TX sticks, observe channel values in FC software
  \item Roll stick (right stick L/R): Channel 1 should move 1000-2000
  \item Pitch stick (right stick U/D): Channel 2 should move 1000-2000
  \item Throttle stick (left stick U/D): Channel 3 should move 1000-2000
  \item Yaw stick (left stick L/R): Channel 4 should move 1000-2000
  \item Screenshot of receiver channels responding
\end{itemize}

WEEK 2 DELIVERABLES (Days 8-14)\par

\begin{itemize}
  \item ESC Configuration \& Motor Test
  \item [WARNING] CRITICAL SAFETY: Remove propeller for ALL indoor testing
  \item Verify propeller is REMOVED
  \item Wire ESC signal cable to FC Motor 1 output:
  \item ESC Signal (yellow/white) $\rightarrow$ FC M1 signal pad
  \item ESC Ground (brown/black) $\rightarrow$ FC GND
  \item ESC +5V (red) $\rightarrow$ DO NOT CONNECT (clip this wire)
  \item In FC software, configure Motor 1 output for motor
  \item Calibrate ESC (follow ESC calibration procedure):
  \item 1. Disconnect battery
  \item 2. Set TX throttle to MAXIMUM
  \item 3. Connect battery (ESC beeps)
  \item 4. Wait for confirmation beeps
  \item 5. Set TX throttle to MINIMUM
  \item 6. ESC beeps to confirm calibration complete
  \item Arm FC (throttle low, arm switch or stick command)
  \item SLOWLY increase throttle to 10-20\%
  \item Motor should spin (NO PROP ATTACHED)
  \item Observe motor direction:
  \item Looking at motor from FRONT of aircraft, motor should spin CLOCKWISE
  \item If COUNTER-CLOCKWISE, swap any two motor wires (ESC to motor connection)
  \item Test throttle response through full range (0-100\%)
  \item Disarm, disconnect battery
  \item Video of motor spin test (upload to shared drive)
\end{itemize}

\begin{itemize}
  \item Servo Wiring \& Configuration
  \item Identify servo outputs on FC: S1, S2, S3, S4 (or PWM1-4)
  \item Wire servos to FC (signal only, power comes from Rail 2):
  \item Servo 1 (Left Aileron):  Signal (orange) $\rightarrow$ FC S1,  Power from Rail 2
  \item Servo 2 (Right Aileron): Signal (orange) $\rightarrow$ FC S2,  Power from Rail 2
  \item Servo 3 (Elevator):      Signal (orange) $\rightarrow$ FC S3,  Power from Rail 2
  \item Servo 4 (Rudder):        Signal (orange) $\rightarrow$ FC S4,  Power from Rail 2
  \item In FC software, configure servo outputs:
  \item Output 1 (S1): Type = Servo, Function = Aileron (or Left Aileron if separate)
  \item Output 2 (S2): Type = Servo, Function = Aileron (or Right Aileron if separate)
  \item Output 3 (S3): Type = Servo, Function = Elevator
  \item Output 4 (S4): Type = Servo, Function = Rudder
  \item For each servo, set:
  \item Min value: 1000µs (may adjust later)
  \item Max value: 2000µs (may adjust later)
  \item Center: 1500µs
  \item Rate: 100\% (adjust in TX if needed)
\end{itemize}

\begin{itemize}
  \item Servo Direction \& Endpoint Configuration
  \item Power system (battery connected, FC + Servos powered)
  \item Verify all 4 servos center when FC armed
  \item Test each control input and verify CORRECT direction:
\end{itemize}

\begin{itemize}
  \item RIGHT STICK LEFT (Roll Left - want left wing down, right wing up):
  \item Servo 1 (Left Aileron): Should move DOWN (trailing edge down)
  \item Servo 2 (Right Aileron): Should move UP (trailing edge up)
  \item If wrong direction, enable REVERSE in FC for that servo
\end{itemize}

\begin{itemize}
  \item RIGHT STICK RIGHT (Roll Right - want right wing down, left wing up):
  \item Servo 1 (Left Aileron): Should move UP
  \item Servo 2 (Right Aileron): Should move DOWN
  \item If wrong direction, enable REVERSE in FC for that servo
\end{itemize}

\begin{itemize}
  \item RIGHT STICK DOWN (Pitch Up - want nose to rise):
  \item Servo 3 (Elevator): Should move UP (trailing edge up)
  \item If wrong direction, enable REVERSE in FC for elevator
\end{itemize}

\begin{itemize}
  \item RIGHT STICK UP (Pitch Down - want nose to lower):
  \item Servo 3 (Elevator): Should move DOWN (trailing edge down)
  \item If wrong direction, enable REVERSE in FC for elevator
\end{itemize}

\begin{itemize}
  \item LEFT STICK LEFT (Yaw Left - want nose to turn left):
  \item Servo 4 (Rudder): Should move LEFT (trailing edge left)
  \item If wrong direction, enable REVERSE in FC for rudder
\end{itemize}

\begin{itemize}
  \item LEFT STICK RIGHT (Yaw Right - want nose to turn right):
  \item Servo 4 (Rudder): Should move RIGHT (trailing edge right)
  \item If wrong direction, enable REVERSE in FC for rudder
\end{itemize}

\begin{itemize}
  \item Document servo directions:
  \item Servo 1: NORMAL / REVERSED
  \item Servo 2: NORMAL / REVERSED
  \item Servo 3: NORMAL / REVERSED
  \item Servo 4: NORMAL / REVERSED
\end{itemize}

\begin{itemize}
  \item Adjust servo endpoints if needed (limit travel to prevent binding):
  \item Manually move control surfaces to check for binding
  \item If binding at full deflection, reduce endpoints in FC (e.g., 1100-1900 instead of 1000-2000)
\end{itemize}

\begin{itemize}
  \item Failsafe Configuration
  \item In FC software, configure failsafe:
  \item Trigger: Loss of radio signal (RX\_LOSS)
\end{itemize}

\subsection{Action}

\begin{itemize}
  \item Throttle $\rightarrow$ 0\% (motor off)
  \item Servos $\rightarrow$ Hold last position OR Center (team decision)
  \item Recommended: Hold position (allows glide)
  \item Duration: Activate after 2 seconds of signal loss
  \item Test failsafe:
  \item 1. Power system, arm FC
  \item 2. Set throttle to 50\%
  \item 3. Turn OFF transmitter
  \item 4. Verify motor cuts within 2 seconds
  \item 5. Verify servos hold position or center
  \item 6. Turn ON transmitter, verify control restored
  \item Failsafe Status: WORKING / NEEDS ADJUSTMENT
\end{itemize}

\begin{itemize}
  \item Documentation
  \item Create wiring diagram showing FC connections (FC $\rightarrow$ GPS, RX, ESC, Servos)
  \item Screenshot of all FC configuration pages (save as PDF)
  \item Write 1-page summary: "Flight Electronics Integration Report"
  \item Firmware version and settings
  \item GPS, receiver, motor, servo configuration
  \item Issues encountered and solutions
  \item Lessons learned
\end{itemize}

\section{SKILLS YOU WILL DEVELOP}

\begin{itemize}
  \item Flight controller firmware configuration (iNav or ArduPilot)
  \item GPS module interfacing and troubleshooting
  \item Radio control system binding and configuration (ELRS)
  \item Servo calibration and direction tuning
  \item ESC calibration and motor direction verification
  \item Failsafe programming for safe flight operations
  \item Systematic testing of control surfaces
\end{itemize}

\section{POTENTIAL ISSUES \& TROUBLESHOOTING}

Problem: GPS won't lock (satellite count = 0)\par

\begin{itemize}
  \item $\rightarrow$ Ensure testing OUTDOORS (GPS signals blocked indoors)
  \item $\rightarrow$ Check GPS wiring (TX/RX may be swapped)
  \item $\rightarrow$ Verify GPS enabled on correct UART in FC software
  \item $\rightarrow$ Try different baud rates (9600, 38400, 115200)
  \item $\rightarrow$ Wait longer (initial lock can take 5-10 minutes)
\end{itemize}

Problem: Receiver not binding to transmitter\par

\begin{itemize}
  \item $\rightarrow$ Ensure RX and TX have same ELRS firmware version (update if mismatch)
  \item $\rightarrow$ Try WiFi binding method instead of button method (or vice versa)
  \item $\rightarrow$ Verify receiver is powered (5V + GND connected)
\end{itemize}

Problem: Receiver channels not showing in FC software\par

\begin{itemize}
  \item $\rightarrow$ Check protocol setting (CRSF vs SBUS)
  \item $\rightarrow$ Verify baud rate matches protocol (420000 for CRSF, 100000 for SBUS)
  \item $\rightarrow$ Check wiring (RX TX $\rightarrow$ FC RX)
  \item $\rightarrow$ Try different UART port on FC
\end{itemize}

Problem: Motor won't spin or spins erratically\par

\begin{itemize}
  \item $\rightarrow$ Verify ESC calibrated correctly
  \item $\rightarrow$ Check throttle channel mapping (usually Channel 3)
  \item $\rightarrow$ Ensure FC armed (check arming flags in software)
  \item $\rightarrow$ Verify motor output configured correctly in FC
\end{itemize}

Problem: Servo moving wrong direction\par

\begin{itemize}
  \item $\rightarrow$ Enable REVERSE in FC for that servo channel
  \item $\rightarrow$ Double-check control surface should move in which direction
\end{itemize}

Problem: Servo jittering or twitching\par

\begin{itemize}
  \item $\rightarrow$ Check power supply (voltage sag? Servos should be on Rail 2)
  \item $\rightarrow$ Reduce servo update rate in FC (try 50Hz instead of 200Hz)
  \item $\rightarrow$ Check for electrical noise (separate signal wires from motor wires)
\end{itemize}

\section{NOTES \& QUESTIONS}

[Space for team to write notes during work]\par

\section{SUBSYSTEM 3: RESEARCH PAYLOAD TEAM}

\section{TEAM MEMBERS}

Owner:         \underline{\hspace{4.8cm}}  Phone: \underline{\hspace{1.7999999999999998cm}}  Email: \underline{\hspace{1.95cm}}\par

Assistant 1:   \underline{\hspace{4.8cm}}  Phone: \underline{\hspace{1.7999999999999998cm}}  Email: \underline{\hspace{1.95cm}}\par

Assistant 2:   \underline{\hspace{4.8cm}}  Phone: \underline{\hspace{1.7999999999999998cm}}  Email: \underline{\hspace{1.95cm}}\par

\section{SUBSYSTEM SCOPE}

You are responsible for the RESEARCH MISSION of the aircraft: collecting environmental\par

data (temperature, humidity, pressure, air quality) and logging it to a MicroSD card.\par

\section{MISSION-CRITICAL IMPORTANCE:}

This is the PURPOSE of the entire project. Without functioning data logging, we have\par

built a drone but accomplished no research. Your work defines success.\par

\section{COMPONENTS ASSIGNED TO YOU}

\begin{itemize}
  \item Arduino Nano (ATmega328P)
  \item HiLetgo BME280 Sensor Module (Temperature, Humidity, Pressure)
  \item Ximimark MQ135 Air Quality Sensor (VOC, CO₂, NH₃)
  \item HiLetgo MicroSD Card Module (SPI interface)
  \item SanDisk 32GB MicroSD Card
  \item Taiss SPST Toggle Switch (2-pin, logging control)
  \item Laptop with Arduino IDE
  \item USB Mini cable (for Arduino programming)
  \item Jumper wires (Dupont M-F or F-F for prototyping)
\end{itemize}

WEEK 1 DELIVERABLES (Days 1-7)\par

\begin{itemize}
  \item Arduino Setup
  \item Install Arduino IDE on laptop (version 1.8.19 or 2.x)
  \item Connect Arduino Nano via USB
  \item Install CH340 drivers if needed (for Nano clones)
  \item Select board: Tools $\rightarrow$ Board $\rightarrow$ Arduino Nano
  \item Select processor: Tools $\rightarrow$ Processor $\rightarrow$ ATmega328P (Old Bootloader if clone)
  \item Select port: Tools $\rightarrow$ Port $\rightarrow$ (select COM port with Nano)
  \item Upload "Blink" example sketch (File $\rightarrow$ Examples $\rightarrow$ 01.Basics $\rightarrow$ Blink)
  \item Verify LED blinks (confirms Arduino functional)
\end{itemize}

\begin{itemize}
  \item BME280 Sensor Setup
  \item Install Adafruit BME280 library:
  \item Sketch $\rightarrow$ Include Library $\rightarrow$ Manage Libraries $\rightarrow$ Search "Adafruit BME280"
  \item Install "Adafruit BME280" library + dependencies
  \item Wire BME280 to Arduino:
  \item BME280 VCC $\rightarrow$ Arduino 5V or 3.3V (check if module is 5V tolerant)
  \item BME280 GND $\rightarrow$ Arduino GND
  \item BME280 SDA $\rightarrow$ Arduino A4 (SDA)
  \item BME280 SCL $\rightarrow$ Arduino A5 (SCL)
  \item Upload I2C Scanner sketch to find BME280 address:
  \item Expected address: 0x76 or 0x77
  \item Actual address found: 0x\_\_\_\_\_\_
  \item Upload BME280 test sketch (from Adafruit library examples)
  \item Open Serial Monitor (9600 baud)
  \item Verify readings:
  \item Temperature: \_\_\_\_\_\_$^\circ$C (should be room temp, \textasciitilde{}20-25$^\circ$C)
  \item Humidity: \_\_\_\_\_\_\% (should be 30-70\% indoors)
  \item Pressure: \_\_\_\_\_\_hPa (should be 980-1020 hPa depending on altitude)
  \item Test repeatability: Readings should be stable, not erratic
  \item Screenshot of Serial Monitor with BME280 readings
\end{itemize}

\begin{itemize}
  \item MQ135 Sensor Setup
  \item Wire MQ135 to Arduino:
  \item MQ135 VCC $\rightarrow$ Arduino 5V (MQ135 requires 5V for heater)
  \item MQ135 GND $\rightarrow$ Arduino GND
  \item MQ135 AOUT (Analog Out) $\rightarrow$ Arduino A0
  \item (DOUT not used in this project)
  \item Upload simple analog read sketch:
  \item int gasValue = analogRead(A0);
  \item Serial.println(gasValue);
  \item Open Serial Monitor (9600 baud)
  \item Verify analog readings:
  \item Initial readings: \_\_\_\_\_\_ (0-1023 range)
  \item Note: MQ135 requires 24-48 hour burn-in for accurate readings
  \item During first test, readings will be unstable - this is normal
  \item Test sensitivity: Blow on sensor, value should change
  \item Document baseline value: \_\_\_\_\_\_ (this is "clean air" reference)
\end{itemize}

\begin{itemize}
  \item MicroSD Card Module Setup
  \item Install SD library (included with Arduino IDE, no need to install separately)
  \item Format MicroSD card as FAT32 (use computer)
  \item Card size <32GB recommended (larger may have compatibility issues)
  \item Insert card into MicroSD module
  \item Wire MicroSD module to Arduino:
  \item SD VCC $\rightarrow$ Arduino 5V
  \item SD GND $\rightarrow$ Arduino GND
  \item SD MISO (DO) $\rightarrow$ Arduino D12 (MISO)
  \item SD MOSI (DI) $\rightarrow$ Arduino D11 (MOSI)
  \item SD SCK (CLK) $\rightarrow$ Arduino D13 (SCK)
  \item SD CS $\rightarrow$ Arduino D10 (or any digital pin, define in code)
  \item Upload SD CardInfo example sketch (File $\rightarrow$ Examples $\rightarrow$ SD $\rightarrow$ CardInfo)
  \item Open Serial Monitor (9600 baud)
  \item Verify SD card detected:
  \item Card type: \underline{\hspace{1.5cm}} (SD, SDHC, or SDXC)
  \item Volume type: \underline{\hspace{1.5cm}} (should be FAT16 or FAT32)
  \item Volume size: \underline{\hspace{1.5cm}} MB
  \item Upload SD ReadWrite example (File $\rightarrow$ Examples $\rightarrow$ SD $\rightarrow$ ReadWrite)
  \item Verify file creation:
  \item Remove SD card, insert into computer
  \item Check for "test.txt" file
  \item Open file, verify text written correctly
  \item File creation Status: SUCCESS / FAILED
\end{itemize}

\begin{itemize}
  \item Toggle Switch Setup
  \item Wire toggle switch to Arduino:
  \item Switch Pin 1 $\rightarrow$ Arduino D2
  \item Switch Pin 2 $\rightarrow$ Arduino GND
  \item Upload sketch to read switch state:
  \item pinMode(2, INPUT\_PULLUP);
  \item int switchState = digitalRead(2);
  \item Serial.println(switchState);
  \item Open Serial Monitor
  \item Test switch:
  \item Switch OPEN: Serial Monitor shows 1 (HIGH)
  \item Switch CLOSED: Serial Monitor shows 0 (LOW)
  \item Switch Status: WORKING / NOT WORKING
\end{itemize}

WEEK 2 DELIVERABLES (Days 8-14)\par

\begin{itemize}
  \item Integrated Data Logging Code Development
  \item Write Arduino sketch that combines all sensors + SD logging:
\end{itemize}

\subsection{Required functionality}

\begin{itemize}
  \item 1. Initialize BME280, MQ135, SD card, toggle switch
  \item 2. Check if SD card is present and writable
  \item 3. Create CSV file with headers:
  \item "timestamp,temp\_C,humidity\_\%,pressure\_hPa,gas\_level"
  \item 4. Read toggle switch state (D2)
  \item 5. If switch CLOSED (LOW):
  \item Read BME280 (temp, humidity, pressure)
  \item Read MQ135 (analog value from A0)
  \item Create timestamp (milliseconds since startup)
  \item Write CSV line to SD card
  \item Flush SD buffer (ensure data written immediately)
  \item Optional: Blink LED to indicate logging
  \item 6. If switch OPEN (HIGH):
  \item Do not log (or log at reduced rate)
  \item 7. Repeat every 1-5 seconds (team decision on logging rate)
\end{itemize}

\begin{itemize}
  \item Code structure:
  \item Use functions for readability:
  \item void readBME280() \{ ... \}
  \item void readMQ135() \{ ... \}
  \item void logToSD() \{ ... \}
  \item Add comments explaining each section
  \item Include error handling (SD card removed, sensor fail, etc.)
\end{itemize}

\begin{itemize}
  \item Upload code to Arduino
  \item Open Serial Monitor to verify readings printing correctly
\end{itemize}

\begin{itemize}
  \item Bench Testing \& Validation
  \item Test 1: SD Card Detection
  \item Power Arduino with SD card inserted
  \item Verify Serial Monitor shows "SD card initialized"
  \item If fails, check wiring and card format
\end{itemize}

\begin{itemize}
  \item Test 2: Sensor Data Validity
  \item Let MQ135 warm up for 5 minutes
  \item Verify BME280 readings reasonable:
  \item Temp: 15-30$^\circ$C (room temp)
  \item Humidity: 20-80\% (typical indoor range)
  \item Pressure: 980-1020 hPa (varies by altitude and weather)
  \item Verify MQ135 analog value: 0-1023 (any value in range is okay for now)
\end{itemize}

\begin{itemize}
  \item Test 3: Toggle Switch Control
  \item Set switch to CLOSED (logging enabled)
  \item Verify Serial Monitor shows logging messages every 1-5 seconds
  \item Set switch to OPEN (logging disabled)
  \item Verify logging stops (or reduced rate)
  \item Set switch CLOSED again
  \item Verify logging resumes
\end{itemize}

\begin{itemize}
  \item Test 4: CSV File Creation
  \item Run logging for 1 minute with switch CLOSED
  \item Power off Arduino
  \item Remove SD card, insert into computer
  \item Verify CSV file exists:
  \item File name: \underline{\hspace{1.5cm}} (e.g., "datalog.csv" or timestamped name)
  \item File size: \underline{\hspace{1.5cm}} bytes (should be >0)
  \item Open CSV file in Excel or text editor
  \item Verify columns:
  \item timestamp | temp\_C | humidity\_\% | pressure\_hPa | gas\_level
  \item \_\_\_\_\_\_\_\_|\_\_\_\_\_\_\_\_|\underline{\hspace{1.7999999999999998cm}}|\underline{\hspace{2.1cm}}|\underline{\hspace{1.5cm}}
  \item \_\_\_\_\_\_\_\_|\_\_\_\_\_\_\_\_|\underline{\hspace{1.7999999999999998cm}}|\underline{\hspace{2.1cm}}|\underline{\hspace{1.5cm}}
  \item \_\_\_\_\_\_\_\_|\_\_\_\_\_\_\_\_|\underline{\hspace{1.7999999999999998cm}}|\underline{\hspace{2.1cm}}|\underline{\hspace{1.5cm}}
  \item Check data is reasonable (no NaN, no -999 errors)
  \item File Status: VALID / ERRORS FOUND
\end{itemize}

\begin{itemize}
  \item MQ135 Burn-In Procedure (CRITICAL)
  \item [WARNING] WARNING: MQ135 requires 24-48 hours of continuous power before accurate readings
\end{itemize}

\begin{itemize}
  \item Plan burn-in procedure:
  \item Option A: Power Arduino + MQ135 continuously for 48 hours before first flight
  \item Connect Arduino to USB power (not battery yet)
  \item Let run for 48 hours in ventilated area
  \item Option B: Power MQ135 separately from bench power supply
  \item 5V to VCC, GND to GND
  \item Leave powered for 48 hours
\end{itemize}

\begin{itemize}
  \item Burn-in start date/time: \underline{\hspace{4.05cm}}
  \item Expected completion: \underline{\hspace{4.05cm}}
  \item During burn-in, monitor sensor readings every 6-12 hours:
  \item Hour 0: Gas value = \_\_\_\_\_\_
  \item Hour 6: Gas value = \_\_\_\_\_\_
  \item Hour 12: Gas value = \_\_\_\_\_\_
  \item Hour 24: Gas value = \_\_\_\_\_\_
  \item Hour 48: Gas value = \_\_\_\_\_\_ (should stabilize by now)
  \item Burn-in Status: COMPLETE / IN PROGRESS / NOT STARTED
\end{itemize}

\begin{itemize}
  \item Arduino Power Integration
  \item Determine Arduino power source:
  \item Option A: Power from Rail 1 (Flight-Critical) via 5V + GND
  \item Option B: Power from Rail 2 (Servo Power) via 5V + GND
  \item Team Decision: Arduino powered from RAIL \_\_\_
\end{itemize}

\begin{itemize}
  \item Disconnect Arduino from USB (will power from 5V rail now)
  \item Connect Arduino 5V pin to chosen rail
  \item Connect Arduino GND to common ground bus
  \item Power system (battery connected)
  \item Verify Arduino powers on (LED should light)
  \item Verify logging works (check SD card after test)
\end{itemize}

\begin{itemize}
  \item Final Integration Test (5-Minute Flight Simulation)
  \item Full system powered (battery, FC, servos, Arduino)
  \item Insert fresh SD card (formatted, empty)
  \item Set toggle switch to CLOSED (logging enabled)
  \item Power on, let system run for 5 minutes
  \item Optional: Move sensors around, blow on MQ135, heat BME280 to create data variation
  \item Set toggle switch to OPEN (logging disabled)
  \item Power off
  \item Remove SD card, check CSV file:
  \item Number of data points: \_\_\_\_\_\_ (expect \textasciitilde{}60-300 depending on logging rate)
  \item File size: \_\_\_\_\_\_ bytes
  \item Data looks reasonable? YES / NO
  \item Plot data in Excel (time vs temp, time vs humidity, time vs gas level)
  \item Screenshot of data plot
\end{itemize}

\begin{itemize}
  \item Documentation
  \item Upload Arduino code to shared drive (with comments!)
  \item Create wiring diagram for Arduino + sensors
  \item Write CSV file format specification (column definitions)
  \item Write 1-page summary: "Research Payload Integration Report"
  \item Sensor specifications and ranges
  \item Data logging format
  \item Calibration procedure (if any)
  \item Known limitations (MQ135 not quantitatively calibrated, etc.)
  \item Lessons learned
\end{itemize}

\section{SKILLS YOU WILL DEVELOP}

\begin{itemize}
  \item Arduino programming (C/C++)
  \item I2C communication protocol (BME280)
  \item Analog-to-digital conversion (ADC) reading (MQ135)
  \item SPI communication protocol (MicroSD)
  \item File I/O and CSV formatting
  \item Sensor interfacing and troubleshooting
  \item Data validation and quality control
\end{itemize}

\section{POTENTIAL ISSUES \& TROUBLESHOOTING}

Problem: BME280 not detected (I2C scan shows no device)\par

\begin{itemize}
  \item $\rightarrow$ Check wiring (SDA $\rightarrow$ A4, SCL $\rightarrow$ A5)
  \item $\rightarrow$ Verify power (VCC $\rightarrow$ 5V or 3.3V, check module specs)
  \item $\rightarrow$ Try both I2C addresses (0x76 and 0x77)
  \item $\rightarrow$ Ensure I2C pullup resistors present (usually on breakout module)
\end{itemize}

Problem: MQ135 readings always 0 or always 1023\par

\begin{itemize}
  \item $\rightarrow$ Check wiring (AOUT $\rightarrow$ A0)
  \item $\rightarrow$ Verify 5V power (MQ135 heater requires 5V)
  \item $\rightarrow$ Sensor may be defective (test with multimeter: should read 0-5V on AOUT)
\end{itemize}

Problem: SD card not detected\par

\begin{itemize}
  \item $\rightarrow$ Verify card formatted as FAT32 (not exFAT or NTFS)
  \item $\rightarrow$ Try different SD card (<32GB recommended)
  \item $\rightarrow$ Check wiring (MISO, MOSI, SCK, CS pins)
  \item $\rightarrow$ Ensure VCC connected to 5V (some modules have onboard regulator)
  \item $\rightarrow$ Try different CS pin (change in code: const int chipSelect = 10;)
\end{itemize}

Problem: CSV file empty or corrupted\par

\begin{itemize}
  \item $\rightarrow$ Add SD flush after each write: file.flush();
  \item $\rightarrow$ Ensure file.close() called before power off
  \item $\rightarrow$ Check SD card not full
  \item $\rightarrow$ Verify file opened successfully: if (!file) \{ Serial.println("Error"); \}
\end{itemize}

Problem: Data looks unrealistic (e.g., temp = -999)\par

\begin{itemize}
  \item $\rightarrow$ Sensor not initializing properly (check begin() function return value)
  \item $\rightarrow$ Check sensor wiring and power
  \item $\rightarrow$ Add error handling in code:
  \item if (isnan(temp)) \{ temp = -999; \} // Flag as error
\end{itemize}

\section{NOTES \& QUESTIONS}

[Space for team to write notes during work]\par

\section{SUBSYSTEM 4: INTEGRATION \& TESTING TEAM}

\section{TEAM MEMBERS}

Owner:         \underline{\hspace{4.8cm}}  Phone: \underline{\hspace{1.7999999999999998cm}}  Email: \underline{\hspace{1.95cm}}\par

Assistant 1:   \underline{\hspace{4.8cm}}  Phone: \underline{\hspace{1.7999999999999998cm}}  Email: \underline{\hspace{1.95cm}}\par

Assistant 2:   \underline{\hspace{4.8cm}}  Phone: \underline{\hspace{1.7999999999999998cm}}  Email: \underline{\hspace{1.95cm}}\par

\section{SUBSYSTEM SCOPE}

You are responsible for SYSTEMS INTEGRATION and VALIDATION. Your job is to ensure all\par

subsystems work together, identify issues before flight, and verify the aircraft is\par

electrically ready for operation.\par

\section{MISSION-CRITICAL IMPORTANCE:}

You are the last line of defense before flight. If you miss a faulty connection or\par

electrical issue, we risk losing the aircraft. Attention to detail is everything.\par

\section{COMPONENTS ASSIGNED TO YOU}

\begin{itemize}
  \item Multimeter (voltage, current, continuity, resistance)
  \item Clamp meter (for measuring high currents) [if available]
  \item Inline current sensor (if available)
  \item Wire strippers, cutters, crimpers
  \item Zip ties, foam tape, velcro straps
  \item Heat shrink tubing (various sizes)
  \item Electrical tape
  \item Labels / label maker
  \item Notebook for recording measurements
\end{itemize}

WEEK 1 DELIVERABLES (Days 1-7)\par

\begin{itemize}
  \item Multimeter Training \& Practice
  \item Learn multimeter functions:
  \item DC Voltage (DCV): For measuring battery, BEC outputs (0-20V range)
  \item DC Current (DCA): For measuring current draw (200mA, 2A, 10A ranges)
  \item Continuity (⏚): For checking ground connections (beeps when 0$\Omega$)
  \item Resistance ($\Omega$): For checking short circuits or open connections
\end{itemize}

\begin{itemize}
  \item Practice measurements on bench:
  \item Test 1: Measure AA battery voltage (expect \textasciitilde{}1.5V)
  \item Measured: \_\_\_\_\_\_V
  \item Test 2: Measure 9V battery voltage (expect \textasciitilde{}9V)
  \item Measured: \_\_\_\_\_\_V
  \item Test 3: Check continuity of wire (should beep)
  \item Status: PASS / FAIL
  \item Test 4: Measure resistance of wire (expect <1$\Omega$)
  \item Measured: \_\_\_\_\_\_$\Omega$
\end{itemize}

\begin{itemize}
  \item Create multimeter quick reference guide (printout for team):
  \item How to measure voltage (red to +, black to -)
  \item How to measure current (in series with circuit)
  \item How to check continuity (across connection)
  \item Safety: Never measure current in voltage mode (blows fuse!)
\end{itemize}

\begin{itemize}
  \item Component Inventory \& Condition Check
  \item Work with all subsystem teams to verify component condition
  \item For each major component, check:
  \item Physical damage (cracks, bent pins, damaged wires)
  \item Connector condition (secure, not loose or damaged)
  \item Wire insulation (no exposed copper, no fraying)
\end{itemize}

\begin{itemize}
  \item Document issues found:
  \item Component: \underline{\hspace{2.4cm}}  Issue: \underline{\hspace{2.4cm}}  Resolution: \underline{\hspace{2.4cm}}
  \item Component: \underline{\hspace{2.4cm}}  Issue: \underline{\hspace{2.4cm}}  Resolution: \underline{\hspace{2.4cm}}
\end{itemize}

\begin{itemize}
  \item Common Ground Topology Planning
  \item Study system architecture (refer to Block Diagram document)
  \item Identify all ground connections needed:
  \item Battery negative (master ground)
  \item ESC ground
  \item External UBEC ground
  \item Flight Controller ground
  \item Arduino ground
  \item Servo grounds (4$\times$)
  \item Sensor grounds (BME280, MQ135, MicroSD)
\end{itemize}

\begin{itemize}
  \item Create ground wiring plan:
\end{itemize}

\subsection{Use STAR TOPOLOGY from battery negative}

\begin{itemize}
  \item Battery (-) $\rightarrow$ Main GND bus $\rightarrow$ ESC, UBEC, FC, Arduino, Servos, Sensors
\end{itemize}

\begin{itemize}
  \item Calculate total wire length needed for ground bus: \_\_\_\_\_\_ meters
  \item Recommended wire gauge for ground: 18-20 AWG (thick enough for current return)
\end{itemize}

\begin{itemize}
  \item Wire Routing \& Management Planning
  \item Study airframe layout (coordinate with Mechanical team)
  \item Plan wire routing:
  \item High-current wires (motor, battery): Shortest path, away from sensitive electronics
  \item Signal wires (servo signals, GPS, RX): Away from motor wires (noise reduction)
  \item Sensor wires: Away from motor/ESC, accessible for troubleshooting
\end{itemize}

\begin{itemize}
  \item Create wire routing diagram (sketch on paper or CAD):
  \item Mark wire paths in airframe
  \item Identify zip tie attachment points
  \item Plan for strain relief at connectors
  \item Allow slack for vibration (don't pull wires tight)
\end{itemize}

\begin{itemize}
  \item Estimate wire lengths needed:
  \item Battery to ESC: \_\_\_\_\_\_ cm
  \item Battery to UBEC: \_\_\_\_\_\_ cm
  \item ESC to Motor: \_\_\_\_\_\_ cm
  \item ESC BEC to FC: \_\_\_\_\_\_ cm
  \item UBEC to Servo Rail: \_\_\_\_\_\_ cm
  \item FC to GPS: \_\_\_\_\_\_ cm
  \item FC to RX: \_\_\_\_\_\_ cm
  \item FC to Servos (4$\times$): \_\_\_\_\_\_ cm each
  \item Arduino to Sensors: \_\_\_\_\_\_ cm
\end{itemize}

WEEK 2 DELIVERABLES (Days 8-14)\par

\begin{itemize}
  \item Continuity Testing (CRITICAL - Do this before first power-on)
  \item Set multimeter to continuity mode (⏚ symbol)
  \item Verify all grounds connected (should beep for 0$\Omega$ connections):
\end{itemize}

\begin{itemize}
  \item Test 1: Battery (-) to ESC ground
  \item Measured: \_\_\_\_\_\_$\Omega$    Status: PASS (<1$\Omega$) / FAIL
\end{itemize}

\begin{itemize}
  \item Test 2: Battery (-) to External UBEC ground
  \item Measured: \_\_\_\_\_\_$\Omega$    Status: PASS / FAIL
\end{itemize}

\begin{itemize}
  \item Test 3: Battery (-) to Flight Controller GND pad
  \item Measured: \_\_\_\_\_\_$\Omega$    Status: PASS / FAIL
\end{itemize}

\begin{itemize}
  \item Test 4: Battery (-) to Arduino GND pin
  \item Measured: \_\_\_\_\_\_$\Omega$    Status: PASS / FAIL
\end{itemize}

\begin{itemize}
  \item Test 5: Battery (-) to Servo 1 brown wire (ground)
  \item Measured: \_\_\_\_\_\_$\Omega$    Status: PASS / FAIL
\end{itemize}

\begin{itemize}
  \item Test 6: Battery (-) to BME280 GND
  \item Measured: \_\_\_\_\_\_$\Omega$    Status: PASS / FAIL
\end{itemize}

\begin{itemize}
  \item If ANY test fails (no beep, high resistance):
  \item $\rightarrow$ STOP and trace wiring
  \item $\rightarrow$ Fix ground connection before proceeding
  \item $\rightarrow$ Re-test until all grounds show continuity
\end{itemize}

\begin{itemize}
  \item Voltage Measurements (Validate Power Distribution)
  \item Power system with battery (no motor/servos connected initially)
  \item Measure voltages at key points:
\end{itemize}

\subsection{Battery voltage (XT60 terminals)}

\begin{itemize}
  \item Expected: 14.0-16.8V    Measured: \_\_\_\_\_\_V    Status: PASS / FAIL
\end{itemize}

\subsection{Rail 1 voltage (ESC BEC output, measure at FC 5V pad)}

\begin{itemize}
  \item Expected: 5.0 $\pm$ 0.1V    Measured: \_\_\_\_\_\_V    Status: PASS / FAIL
\end{itemize}

\subsection{Rail 2 voltage (External UBEC output, measure at servo power rail)}

\begin{itemize}
  \item Expected: 5.0 $\pm$ 0.1V    Measured: \_\_\_\_\_\_V    Status: PASS / FAIL
\end{itemize}

\begin{itemize}
  \item GPS module VCC (at GPS connector):
  \item Expected: 5.0 $\pm$ 0.1V    Measured: \_\_\_\_\_\_V    Status: PASS / FAIL
\end{itemize}

\subsection{Receiver VCC}

\begin{itemize}
  \item Expected: 5.0 $\pm$ 0.1V    Measured: \_\_\_\_\_\_V    Status: PASS / FAIL
\end{itemize}

\subsection{Arduino 5V pin}

\begin{itemize}
  \item Expected: 5.0 $\pm$ 0.1V    Measured: \_\_\_\_\_\_V    Status: PASS / FAIL
\end{itemize}

\begin{itemize}
  \item If any voltage out of range:
  \item $\rightarrow$ Check connections (cold solder joint? Loose connector?)
  \item $\rightarrow$ Measure voltage at BEC output vs at component (voltage drop indicates bad connection)
\end{itemize}

\begin{itemize}
  \item Current Measurements (Validate Power Budget)
  \item Measure current draw for each subsystem:
\end{itemize}

\begin{itemize}
  \item Method: Use clamp meter (around wire) or inline current sensor
\end{itemize}

\subsection{Rail 1 idle (only FC + GPS + RX powered, no Arduino yet)}

\begin{itemize}
  \item Expected: \textasciitilde{}320-400mA    Measured: \_\_\_\_\_\_mA    Status: PASS / FAIL
\end{itemize}

\subsection{Rail 1 with Arduino active (logging)}

\begin{itemize}
  \item Expected: \textasciitilde{}660mA    Measured: \_\_\_\_\_\_mA    Status: PASS / FAIL
\end{itemize}

\subsection{Rail 2 idle (servos idle, no movement)}

\begin{itemize}
  \item Expected: \textasciitilde{}40mA    Measured: \_\_\_\_\_\_mA    Status: PASS / FAIL
\end{itemize}

\subsection{Rail 2 with 2 servos moving}

\begin{itemize}
  \item Expected: \textasciitilde{}400-500mA    Measured: \_\_\_\_\_\_mA    Status: PASS / FAIL
\end{itemize}

\subsection{Rail 2 with all 4 servos moving}

\begin{itemize}
  \item Expected: \textasciitilde{}720-1200mA    Measured: \_\_\_\_\_\_mA    Status: PASS / FAIL
\end{itemize}

\subsection{Battery total current (motor at 50\% throttle, NO PROP)}

\begin{itemize}
  \item Expected: \textasciitilde{}10-12A    Measured: \_\_\_\_\_\_A    Status: PASS / FAIL
\end{itemize}

\begin{itemize}
  \item Compare measured currents to Power Budget document predictions
  \item If currents significantly higher than expected (>20\% over):
  \item $\rightarrow$ Investigate (component drawing excessive current? Short circuit?)
\end{itemize}

\begin{itemize}
  \item Voltage Sag Test (MOST CRITICAL TEST)
  \item [WARNING] This test verifies FC won't brown out during servo load
\end{itemize}

\begin{itemize}
  \item Setup:
  \item Connect multimeter probes to Flight Controller 5V input (leave connected)
  \item Power full system (battery connected, all subsystems active)
  \item Set multimeter to DC voltage, 20V range
\end{itemize}

\begin{itemize}
  \item Baseline voltage (idle, no servo movement):
  \item Measured: \_\_\_\_\_\_V    (should be \textasciitilde{}5.0V)
\end{itemize}

\begin{itemize}
  \item Voltage under load (move all 4 servos simultaneously, full deflection):
  \item Measured: \_\_\_\_\_\_V    (MUST stay >4.8V, ideally >4.9V)
\end{itemize}

\begin{itemize}
  \item PASS CRITERIA: Voltage must stay >4.8V during aggressive servo movement
  \item Status: PASS / FAIL
\end{itemize}

\begin{itemize}
  \item If FAIL (voltage drops below 4.8V):
\end{itemize}

\subsection{Possible causes}

\begin{itemize}
  \item Servos powered from Rail 1 instead of Rail 2 (check wiring!)
  \item Excessive current draw (measure with clamp meter)
  \item Insufficient BEC capacity (need larger UBEC?)
  \item High-resistance connection (voltage drop in wiring)
\end{itemize}

\begin{itemize}
  \item Resolution: \underline{\hspace{8.85cm}}
\end{itemize}

\begin{itemize}
  \item GPS Lock Test with Motor Running (Electrical Noise Test)
  \item [WARNING] This test verifies motor/ESC noise doesn't interfere with GPS
\end{itemize}

\begin{itemize}
  \item Setup:
  \item Power full system
  \item Wait for GPS lock (satellite count >8)
  \item Open FC configuration software, monitor GPS satellite count
\end{itemize}

\begin{itemize}
  \item Baseline GPS lock (motor off):
  \item Satellite count: \_\_\_\_\_\_    (should be >8)
\end{itemize}

\begin{itemize}
  \item GPS lock with motor running (50\% throttle, NO PROP, outdoors):
  \item Satellite count: \_\_\_\_\_\_    (should NOT drop significantly)
\end{itemize}

\begin{itemize}
  \item PASS CRITERIA: GPS maintains lock, satellite count drops by <3
  \item Status: PASS / FAIL
\end{itemize}

\begin{itemize}
  \item If FAIL (GPS loses lock or satellite count drops significantly):
\end{itemize}

\subsection{Possible causes}

\begin{itemize}
  \item Motor wires too close to GPS antenna
  \item Motor wires not twisted (increases EMI)
  \item GPS not shielded from motor/ESC
\end{itemize}

\subsection{Mitigation}

\begin{itemize}
  \item Twist motor wires together (reduces electromagnetic interference)
  \item Increase separation between motor/ESC and GPS (>10cm minimum)
  \item Add ferrite bead to motor wires (reduces high-frequency noise)
  \item Shield GPS module with aluminum foil or copper tape (if desperate)
\end{itemize}

\begin{itemize}
  \item Resolution: \underline{\hspace{8.85cm}}
\end{itemize}

\begin{itemize}
  \item Wire Management \& Vibration Isolation
  \item Secure all wires with zip ties (every 10-15cm)
  \item Ensure no wires rub against:
  \item Control surfaces (ailerons, elevator, rudder)
  \item Sharp edges (can cut through insulation)
  \item Moving parts (servo arms, linkages)
  \item Add strain relief at connectors:
  \item Loop wire before connector, secure with zip tie
  \item Prevents connector pull-out during vibration
  \item Mount vibration-sensitive components on foam:
  \item Flight Controller (double-sided foam tape)
  \item GPS module (foam pad)
  \item Arduino + sensors (foam pad or rubber standoffs)
  \item Label all major wires:
  \item Battery + / Battery -
  \item Rail 1 +5V / Rail 1 GND
  \item Rail 2 +5V / Rail 2 GND
  \item Motor wires (Phase A, B, C)
  \item Servo signals (S1, S2, S3, S4)
\end{itemize}

\begin{itemize}
  \item Pre-Flight Electrical Checklist Creation
  \item Create laminated checklist for pre-flight verification:
\end{itemize}

\section{PRE-FLIGHT ELECTRICAL CHECKLIST:}

\begin{itemize}
  \item Battery voltage >14.0V (measure with multimeter)
  \item Battery balance: All cells within 0.05V (check via balance connector)
  \item All connectors secure (XT60, servo connectors, FC connections)
  \item No loose wires or exposed metal
  \item Propeller secure (spinner nut tight) [only when ready for flight]
  \item GPS module LED blinking (indicates searching for lock)
  \item FC status LED normal pattern (not rapid blinking = error)
  \item Receiver LED solid (indicates bound to TX)
  \item Arduino logging active (if switch ON, LED indicates activity)
  \item Battery charged within last 24 hours (LiPo self-discharge)
  \item Fire extinguisher and first aid kit present
  \item LiPo safety bag available (for post-flight storage)
\end{itemize}

\begin{itemize}
  \item Print 5 copies of checklist (use for each flight/test)
\end{itemize}

\begin{itemize}
  \item Full System Integration Test
  \item [WARNING] This is the final validation before airframe installation
\end{itemize}

\begin{itemize}
  \item Setup: All electrical systems on bench, fully wired, battery connected
  \item Power on sequence:
  \item 1. Verify all connections secure
  \item 2. Turn ON transmitter (throttle at zero)
  \item 3. Connect battery
  \item 4. Wait 30 seconds for system initialization
  \item 5. Verify GPS lock (wait 2-5 minutes outdoors if needed)
\end{itemize}

\begin{itemize}
  \item Functional tests:
  \item Flight Controller status LED: Normal pattern? YES / NO
  \item GPS lock: Satellite count >8? \_\_\_\_\_\_ satellites
  \item Receiver control: Move TX sticks, servos respond? YES / NO
  \item Motor response: Arm FC, increase throttle (NO PROP), motor spins? YES / NO
  \item Servo directions: All 4 servos move correct direction? YES / NO
  \item Failsafe: Turn OFF TX, motor cuts within 2 sec? YES / NO
  \item Arduino logging: Toggle switch ON, data logging to SD card? YES / NO
\end{itemize}

\begin{itemize}
  \item Power measurements under full load:
  \item Battery voltage (under load): \_\_\_\_\_\_V (should stay >14.0V)
  \item Battery current draw: \_\_\_\_\_\_A (motor at 50\% throttle)
  \item Flight time estimate: Battery capacity (2.2Ah) / Avg current (A) = \_\_\_\_\_\_ hours
  \item (Multiply by 60 for minutes, expect \textasciitilde{}9-13 minutes)
\end{itemize}

\begin{itemize}
  \item Integration Status: READY FOR AIRFRAME INSTALL / ISSUES FOUND
\end{itemize}

\begin{itemize}
  \item If issues found, document and resolve before proceeding:
  \item Issue 1: \underline{\hspace{9.45cm}}
  \item Issue 2: \underline{\hspace{9.45cm}}
\end{itemize}

\begin{itemize}
  \item Documentation
  \item Take photos of final wiring setup (all angles)
  \item Create "as-built" wiring diagram (showing actual wire colors and routing)
  \item Record all voltage and current measurements in spreadsheet
  \item Write 2-page summary: "System Integration Test Report"
  \item Test procedures and results
  \item Pass/fail status for each test
  \item Issues found and resolutions
  \item Readiness assessment for airframe integration
  \item Recommendations for flight testing
\end{itemize}

\section{SKILLS YOU WILL DEVELOP}

\begin{itemize}
  \item Multimeter proficiency (voltage, current, continuity, resistance)
  \item System-level thinking (understanding interactions between subsystems)
  \item Electrical troubleshooting methodology
  \item Wire management and professional routing
  \item Vibration isolation for electronics
  \item Test procedure development and execution
  \item Documentation and reporting
\end{itemize}

\section{POTENTIAL ISSUES \& TROUBLESHOOTING}

Problem: No continuity between grounds (multimeter doesn't beep)\par

\begin{itemize}
  \item $\rightarrow$ Trace ground wiring, find break or missing connection
  \item $\rightarrow$ Check solder joints (cold joint may have high resistance)
  \item $\rightarrow$ Verify all components share common ground bus
\end{itemize}

Problem: Voltage correct at BEC output but low at component\par

\begin{itemize}
  \item $\rightarrow$ Voltage drop indicates high resistance in wiring
  \item $\rightarrow$ Check wire gauge (may be too thin for current draw)
  \item $\rightarrow$ Check connectors (loose or dirty contacts)
  \item $\rightarrow$ Improve connection (thicker wire, better solder joint)
\end{itemize}

Problem: Current draw much higher than predicted\par

\begin{itemize}
  \item $\rightarrow$ Check for short circuit (measure resistance, should be high when powered off)
  \item $\rightarrow$ Identify which component drawing excessive current (disconnect one at a time)
  \item $\rightarrow$ Component may be defective
\end{itemize}

Problem: Voltage sag test fails (FC voltage drops below 4.8V)\par

\begin{itemize}
  \item $\rightarrow$ CRITICAL ISSUE - Do not fly!
  \item $\rightarrow$ Verify servos powered from Rail 2 (not Rail 1)
  \item $\rightarrow$ Check External UBEC output capacity (may need higher current rating)
  \item $\rightarrow$ Reduce servo load (limit travel in FC to reduce peak current)
\end{itemize}

Problem: GPS loses lock during motor test\par

\begin{itemize}
  \item $\rightarrow$ Electrical noise from motor/ESC interfering with GPS
  \item $\rightarrow$ Twist motor wires together (reduces EMI)
  \item $\rightarrow$ Increase physical separation (>10cm)
  \item $\rightarrow$ Add ferrite bead to motor power wires
  \item $\rightarrow$ Shield GPS module (aluminum foil, copper tape)
\end{itemize}

\section{NOTES \& QUESTIONS}

[Space for team to write notes during work]\par

\section{FINAL REMINDERS FOR ALL TEAMS}

\begin{itemize}
  \item Safety first, always. If you're unsure, ask.
  \item Document everything. Photos, measurements, notes. This is your portfolio.
  \item Test early and often. Don't wait until integration to discover issues.
  \item Communicate with other teams. Your subsystems must work together.
  \item Ask questions. No question is stupid. Confusion now prevents crashes later.
  \item Meet deadlines. Other teams depend on your deliverables.
  \item Take pride in your work. This is aerospace engineering, not a hobby project.
\end{itemize}

You have the knowledge, the tools, and the team. Now execute.\par

Build something great.\par

\end{document}
